\chapter{Networking: GNAT.Sockets over the W5500}
\label{ch:networking}

Storage gave the application a filesystem; this chapter gives it a network. The
striking thing about the networking layer is the API the application sees: it is
the \emph{standard} \texttt{GNAT.Sockets} --- \texttt{Create\_Socket},
\texttt{Connect\_Socket}, \texttt{Send\_Socket}, a \texttt{Stream} over a socket.
Portable networking code compiles and runs here unchanged, within the subset this
chapter describes. Underneath that familiar surface sits a small stack: a
chip-neutral interface every NIC implements, and the WIZnet W5500 driver as the
first chip behind it.

\begin{figure}[ht]
\centering
\begin{tikzpicture}[node distance=2mm]
\node[layer, fill=blue!9] (app)  {Your application --- portable \texttt{GNAT.Sockets} code, \texttt{DNS\_Client}};
\node[layer, below=of app]  (fac)  {\texttt{GNAT.Sockets} facade --- interface registry, address types, the \texttt{Stream}};
\node[layer, below=of fac]  (intf) {\texttt{Net\_Devices.Device'Class} --- the chip-neutral contract (dispatched on)};
\node[layer, below=of intf] (wrap) {\texttt{ESP32S3.W5500.Net\_Device} --- the W5500 as the first backend};
\node[layer, below=of wrap] (eng)  {\texttt{ESP32S3.W5500.Sockets} --- the socket engine (\texttt{.Interrupts}, \texttt{.DHCP} alongside)};
\node[layer, below=of eng]  (xport){\texttt{ESP32S3.W5500} --- VDM frame transport + bring-up};
\node[layer, below=of xport] (spi) {\texttt{ESP32S3.SPI} --- task-safe SPI master $\rightarrow$ the chip};
\end{tikzpicture}
\caption{The networking stack. Everything above \texttt{Net\_Devices.Device'Class}
is chip-independent; everything below it is the W5500 driver. A second NIC slots in
as another backend without the layers above noticing.}
\label{fig:net-layers}
\end{figure}

\cnote{On an ESP-IDF system you reach for lwIP and the BSD socket API, with the
Wi-Fi or Ethernet \texttt{netif} underneath. The shape here is the same --- a
sockets API over a driver --- but the W5500 carries its \emph{own} TCP/IP stack in
silicon, so there is no lwIP to link: the ``stack'' is eight hardware sockets on
the chip, and our job is only to drive them and present the sockets API.}

\section{The W5500 and its driver}\index{W5500}\index{Ethernet}\index{SPI}
\label{ch:net:w5500}

The WIZnet W5500 is an Ethernet controller with a \emph{hardwired} TCP/IP stack:
ARP, ICMP, IP, TCP and UDP are implemented in the chip, which exposes eight
independent \emph{hardware sockets} and 32\,KB of buffer RAM, and is driven over
SPI. You do not run a protocol stack on the host; you write a destination address
and payload into a socket's registers and tell it to \texttt{SEND}. On this board
the chip is wired to SPI2 with \texttt{MISO=IO45}, \texttt{MOSI=IO4},
\texttt{SCLK=IO1}, \texttt{SCSn=IO39}, \texttt{INTn=IO3} (pull-up) and
\texttt{RSTn=IO11} (pull-up); the chip-select is an ordinary GPIO held low across
each frame, exactly as the SD-card driver does.

The driver is built in layers, each a child package, so the application can stop at
whichever level it needs:

\begin{itemize}[nosep]
\item \texttt{ESP32S3.W5500} --- the transport. The W5500's SPI ``variable data
length'' frame is a 16-bit address, a control byte (block-select + read/write), and
the data; this layer reads and writes registers and the socket buffers over the
task-safe SPI master, and brings the chip up (reset, \texttt{VERSIONR} probe, MAC /
IP / subnet / gateway, PHY link).
\item \texttt{ESP32S3.W5500.Sockets} --- the \emph{socket engine}. A self-contained
\texttt{Socket} handle with \texttt{Open\_TCP} / \texttt{Open\_UDP},
\texttt{Listen}, \texttt{Connect}, \texttt{Send} / \texttt{Receive},
\texttt{Send\_To} / \texttt{Receive\_From}, and a \texttt{Status} error model. This
is the level the higher layers build on.
\item \texttt{ESP32S3.W5500.Interrupts} --- optional. The single \texttt{INTn} line
(falling edge on \texttt{IO3}) wakes tasks blocked in \texttt{Wait\_Data} /
\texttt{Wait\_Connected} via per-socket suspension objects, so they sleep instead
of poll. Leaving it out costs nothing --- the engine polls.
\item \texttt{ESP32S3.W5500.DHCP} --- a DHCP client (see
\S\ref{ch:net:dns-dhcp}).
\item \texttt{ESP32S3.W5500.Net\_Device} --- the adapter that presents the chip as a
\texttt{GNAT.Sockets} interface (\S\ref{ch:net:netdevice}).
\end{itemize}

The engine and everything above it use controlled SPI handles and
\texttt{Ada.Real\_Time}, so the networking stack requires the \texttt{embedded} or
\texttt{full} profile.

\section{What GNAT.Sockets provides}\index{GNAT.Sockets}\index{socket}
\label{ch:net:gnatsockets}

\texttt{GNAT.Sockets} on a desktop is GNAT's binding to the operating system's
socket API. There is no operating system here, so this is a hand-written subset of
the \emph{same} package, backed by the W5500. Code written to it reads like any
sockets program:

\begin{lstlisting}
type Port_Type is range 0 .. 65535;
type Family_Type is (Family_Inet);                 --  IPv4
type Mode_Type   is (Socket_Stream, Socket_Datagram);

type Inet_Addr_Type is private;
function Inet_Addr (Image : String) return Inet_Addr_Type;   --  "a.b.c.d"
function Image (Value : Inet_Addr_Type) return String;
Any_Inet_Addr : constant Inet_Addr_Type;                     --  0.0.0.0

type Sock_Addr_Type is record
   Family : Family_Type := Family_Inet;
   Addr   : Inet_Addr_Type;
   Port   : Port_Type := 0;
end record;

type Socket_Type is private;
Socket_Error : exception;

procedure Create_Socket  (Socket : out Socket_Type; Family : Family_Type := Family_Inet;
                          Mode : Mode_Type := Socket_Stream);
procedure Bind_Socket    (Socket : in out Socket_Type; Address : Sock_Addr_Type);
procedure Listen_Socket  (Socket : in out Socket_Type; Length : Natural := 15);
procedure Accept_Socket  (Server : Socket_Type; Socket : out Socket_Type;
                          Address : out Sock_Addr_Type);
procedure Connect_Socket (Socket : in out Socket_Type; Server : Sock_Addr_Type);
procedure Send_Socket    (Socket : Socket_Type; Item : Stream_Element_Array;
                          Last : out Stream_Element_Offset; To : access Sock_Addr_Type := null);
procedure Receive_Socket (Socket : Socket_Type; Item : out Stream_Element_Array;
                          Last : out Stream_Element_Offset; From : access Sock_Addr_Type := null);
procedure Close_Socket   (Socket : in out Socket_Type);
function  Stream (Socket : Socket_Type) return Stream_Access;
\end{lstlisting}

A \texttt{Receive\_Timeout} socket option is provided through the standard
\texttt{Set\_Socket\_Option}, so a blocked receive can give up rather than hang:

\begin{lstlisting}
Set_Socket_Option (Sock, Socket_Level, (Receive_Timeout, Timeout => 3.0));
\end{lstlisting}

When the timeout elapses, \texttt{Receive\_Socket} raises \texttt{Socket\_Error},
just as it does on a desktop --- so portable code that already guards against that
works unchanged.

A few necessary departures from desktop \texttt{GNAT.Sockets}:

\begin{itemize}[nosep]
\item \textbf{IPv4 only} (\texttt{Family\_Inet}); the backend has no IPv6 yet.
\item \textbf{One bare-metal init.} Desktop \texttt{Initialize} takes no argument;
here you must say which chip backs the facade (\S\ref{ch:net:netdevice}).
\item \textbf{A finite number of sockets} --- the sum the interfaces provide (eight
per W5500). \texttt{Create\_Socket} past that raises \texttt{Socket\_Error}.
\item \textbf{Accept.} The W5500's listening socket \emph{becomes} the connection,
so \texttt{Accept\_Socket} returns the listener itself as the connected socket; to
accept another client, \texttt{Listen\_Socket} again.
\end{itemize}

\section{Wiring a driver in: the Net\_Device interface}\index{Net\_Devices}
\label{ch:net:netdevice}

The facade does not talk to the W5500 directly. Between them sits one small
package, \texttt{Net\_Devices}, defining the contract a network interface must
satisfy --- the part the facade actually needs. It is an Ada \texttt{interface}
(dispatched on at run time), plus chip-neutral address and status types:

\begin{lstlisting}
type Device is limited interface;
type Device_Access is access all Device'Class;

function Socket_Count (Self : Device) return Positive is abstract;
function Local_IP     (Self : Device) return IPv4_Address is abstract;
function Subnet_Mask  (Self : Device) return IPv4_Address is abstract;

procedure Open (Self : in out Device; Index : Natural; Mode : Transport;
                Local_Port : Port_Number; Result : out Status) is abstract;
procedure Close          (Self : in out Device; Index : Natural) is abstract;
procedure Listen         (Self : in out Device; Index : Natural; Result : out Status) is abstract;
procedure Wait_Connected (Self : in out Device; Index : Natural; Result : out Status) is abstract;
procedure Peer           (Self : in out Device; Index : Natural;
                          Addr : out IPv4_Address; Port : out Port_Number) is abstract;
procedure Connect        (Self : in out Device; Index : Natural; Host : IPv4_Address;
                          Port : Port_Number; Result : out Status) is abstract;
procedure Wait_Data      (Self : in out Device; Index : Natural; Result : out Status) is abstract;
procedure Send           (Self : in out Device; Index : Natural; Data : Stream_Element_Array;
                          Sent : out Natural; Result : out Status) is abstract;
procedure Receive        (Self : in out Device; Index : Natural; Into : out Stream_Element_Array;
                          Count : out Natural; Result : out Status) is abstract;
procedure Send_To        (Self : in out Device; Index : Natural; Host : IPv4_Address;
                          Port : Port_Number; Data : Stream_Element_Array; Result : out Status) is abstract;
procedure Receive_From   (Self : in out Device; Index : Natural; From : out IPv4_Address;
                          From_Port : out Port_Number; Into : out Stream_Element_Array;
                          Count : out Natural; Result : out Status) is abstract;
procedure Set_Receive_Timeout (Self : in out Device; Index : Natural; To : Duration) is abstract;
\end{lstlisting}

Sockets are addressed by an \texttt{Index} the device maps to its own per-socket
state (the W5500's eight hardware sockets). This is the \emph{offloaded-stack}
model: the device hands out TCP/UDP sockets directly. A raw-MAC chip with no stack
of its own would instead drive a software TCP/IP stack that implements this same
interface.

\subsection{The W5500 as a backend}

\texttt{ESP32S3.W5500.Net\_Device} is the first concrete \texttt{Net\_Devices.%
Device}. An \texttt{Instance} wraps one chip and owns that chip's eight
\texttt{Sockets.Socket} handles; each operation delegates to the engine and
converts the (identically-ordered) status and the address bytes. For the common
single-NIC board, one call wires everything together:

\begin{lstlisting}
type Instance is limited new Net_Devices.Device with private;
procedure Attach (Self : in out Instance; Dev : access ESP32S3.W5500.Device);
--  Single-NIC convenience: attach a built-in Instance and register it as default.
procedure Register_Default (Dev : access ESP32S3.W5500.Device);
\end{lstlisting}

The facade keeps a small \emph{registry} of interfaces and dispatches every socket
operation through it; a \texttt{Socket\_Type} records which interface it lives on
and which of that interface's sockets it is. Registration returns an id:

\begin{lstlisting}
type Interface_Id is range 0 .. 3;
function  Add_Interface (Device : Net_Devices.Device_Access) return Interface_Id;
procedure Initialize    (Device : Net_Devices.Device_Access);   --  register as default
\end{lstlisting}

So a board's bring-up ends with \texttt{ESP32S3.W5500.Net\_Device.Register\_Default
(Dev'Access)}, and from there the code is ordinary \texttt{GNAT.Sockets}. Because
the facade names only \texttt{Net\_Devices}, supporting a different chip means
writing one new backend (another implementation of the interface) and registering
it --- nothing above \texttt{Net\_Devices.Device'Class} changes, and a board may
register more than one. \emph{Per-destination routing across interfaces is not
wired yet: today every socket uses the default interface; the registry and the
per-socket interface id are the hooks a routing layer will use.}

\cnote{The interface plays the role of a Linux \texttt{net\_device} /
\texttt{struct net\_device\_ops}, or a \texttt{netif} driver in lwIP: a fixed set
of operations a chip implements so the protocol layer above stays generic. Here the
``protocol layer'' is the W5500's silicon; the interface is just how the facade
reaches it.}

\section{DNS the portable way, DHCP the chip way}\index{DNS}\index{DHCP}
\label{ch:net:dns-dhcp}

Two protocols sit just above sockets, and where they live tells you something about
the architecture.

\textbf{DNS} is an ordinary UDP request/response that happens \emph{after} the
interface is up, so it needs nothing chip-specific. It is written purely against
\texttt{GNAT.Sockets} and lives in a standalone, portable package --- the same
source would resolve names on a desktop:

\begin{lstlisting}
--  in DNS_Client
function Resolve (Server : Inet_Addr_Type; Name : String;
                  Addr : out Inet_Addr_Type;
                  Timeout : Duration := 0.0;
                  Local_Port : Port_Type := 13_001) return Boolean;
\end{lstlisting}

\textbf{DHCP} is the opposite kind of protocol: its job is to \emph{bootstrap} the
interface --- it must talk before an address exists and then program the obtained
IP, mask and gateway into the NIC. Both are operations outside the sockets API on
every platform (raw sockets and \texttt{ioctl} on a desktop; the W5500's
\texttt{Configure} registers here). So the DHCP client cannot be portable the way
DNS is; it rightly lives at the chip layer, \texttt{ESP32S3.W5500.DHCP}, built
straight on the socket engine. It runs the DORA exchange (Discover / Offer /
Request / Ack) and, on success, configures the chip. \texttt{Maintain} starts a
background task that acquires a lease and then keeps it --- renewing at $T1$,
rebinding at $T2$, re-acquiring on expiry:

\begin{lstlisting}
procedure Maintain (Dev : Sockets.Device_Access; MAC : MAC_Address;
                    Socket : Socket_Id := 0; On_Bound : Bound_Callback := null);
\end{lstlisting}

\section{The examples}
\label{ch:net:examples}

Several examples in \texttt{examples/} exercise the stack from the bottom up. Each
shares a tiny \texttt{W5500\_Dev} bring-up package (SPI + reset + a static address,
then \texttt{Register\_Default}); the DHCP example configures itself instead. (FTP,
the most demanding client, gets its own section below.)

\subsection{\texttt{esp32s3\_w5500} --- TCP echo server}\index{TCP}

The full stack end to end: bring-up, the optional interrupts, and the facade. It
listens on \texttt{192.168.1.50:5000} and echoes each line back. This is the TCP
\emph{server} path --- \texttt{Listen\_Socket}, \texttt{Accept\_Socket} (which
reports the peer), \texttt{Receive\_Socket}, \texttt{Send\_Socket} --- with tasks
sleeping on \texttt{INTn} rather than polling.

\begin{lstlisting}
Create_Socket (Server, Family_Inet, Socket_Stream);
Bind_Socket   (Server, (Family_Inet, Any_Inet_Addr, 5000));
loop
   Listen_Socket (Server);
   Accept_Socket (Server, Conn, Peer);          --  blocks until a client connects
   loop
      Receive_Socket (Conn, Buf, Last);
      exit when Last < Buf'First;               --  peer closed
      Send_Socket (Conn, Buf (Buf'First .. Last), Sent);
   end loop;
   Close_Socket (Conn);
end loop;
\end{lstlisting}

\subsection{\texttt{esp32s3\_w5500\_http} --- HTTP GET (TCP client)}

The TCP \emph{client} path. It connects out to a web server, sends
\texttt{GET / HTTP/1.0}, and prints the response until the server closes ---
\texttt{Connect\_Socket} then a \texttt{Receive\_Socket} loop ending at
end-of-stream (\texttt{Last < Item'First}).

\subsection{\texttt{esp32s3\_w5500\_ntp} --- network time (UDP)}\index{NTP}\index{UDP}

A UDP datagram exchange: it sends a 48-byte NTP request to a time server with
\texttt{Send\_Socket (\dots, To => Server'Access)}, reads the reply with
\texttt{Receive\_Socket (\dots, From => \dots)}, takes the transmit timestamp, and
converts the NTP seconds to a civil date/time. Demonstrates \texttt{Socket\_%
Datagram} and the \texttt{To} / \texttt{From} addresses.

\subsection{\texttt{esp32s3\_w5500\_dns} --- name resolution}

A thin demonstration of \texttt{DNS\_Client.Resolve}: it resolves a hostname through
a resolver and prints the address. The protocol --- building the query, parsing the
answer with name-compression --- lives in the reusable module, so the example is a
few lines:

\begin{lstlisting}
if DNS_Client.Resolve (Inet_Addr ("8.8.8.8"), "example.com", Addr, Timeout => 3.0) then
   Put_Line ("example.com = " & Image (Addr));
end if;
\end{lstlisting}

\subsection{\texttt{esp32s3\_w5500\_dhcp} --- a self-configuring address}

No static IP: the board leases one. It calls \texttt{DHCP.Maintain} with a
library-level \texttt{On\_Bound} callback that prints each (re)bind, then the
background task acquires and keeps the lease for the life of the program. After the
first bind the chip is configured, so the socket engine and the facade are ready.

\section{Tying it together: the current weather}
\label{ch:net:weather}

\texttt{esp32s3\_w5500\_weather} uses the whole stack to answer a real question:
\emph{what is the weather at a latitude and longitude?} It fetches the current
conditions from Open-Meteo (\texttt{open-meteo.com}). The flow touches every layer
we have built --- UDP and TCP, name resolution, a receive timeout --- and is worth
reading start to finish.

Open-Meteo's forecast API answers over plain HTTP on port 80. That matters: the
W5500 has no TLS, so an HTTPS-only service would be unreachable without a software
crypto stack. Plain HTTP makes this a straight socket program.

\textbf{1. Bring up the interface.} \texttt{W5500\_Dev.Bring\_Up} does SPI + reset,
sets a static address, waits briefly for the PHY link, and registers the chip:
\texttt{ESP32S3.W5500.Net\_Device.Register\_Default (Dev'Access)}. From here it is
all \texttt{GNAT.Sockets}.

\textbf{2. Resolve the host (UDP).} The server's address is not hard-coded; it is
resolved at run time with the portable resolver, which sends a DNS query over a UDP
socket and parses the reply:

\begin{lstlisting}
if not DNS_Client.Resolve (DNS_Server, "api.open-meteo.com", Server_IP, Timeout => 5.0) then
   Put_Line ("[wx] DNS resolution failed");  return;
end if;
\end{lstlisting}

\textbf{3. Fetch the forecast (TCP).} Open a stream socket, connect to the resolved
address on port 80, and send a one-line HTTP request built from the latitude and
longitude; then read the whole response until the server closes:

\begin{lstlisting}
Create_Socket  (Sock, Family_Inet, Socket_Stream);
Connect_Socket (Sock, (Family_Inet, Server_IP, 80));
Send_Socket    (Sock, To_SEA (Request), Last);
loop
   Receive_Socket (Sock, Buf, Last);
   exit when Last < Buf'First;                  --  server closed the connection
   --  ... append Buf (Buf'First .. Last) to the response text ...
end loop;
Close_Socket (Sock);
\end{lstlisting}

where the request is

\begin{lstlisting}
"GET /v1/forecast?latitude=" & Latitude & "&longitude=" & Longitude
  & "&current_weather=true HTTP/1.0" & CRLF & "Host: api.open-meteo.com" & CRLF
  & "Connection: close" & CRLF & CRLF;
\end{lstlisting}

\textbf{4. Read the answer.} The reply is JSON. A small self-contained scrape (no
heap, no JSON library) pulls the fields out of the \texttt{current\_weather}
object --- anchoring past the \texttt{current\_weather\_units} object first, whose
same-named keys are unit \emph{strings} --- and a table turns the WMO weather code
into a word. The result, for Atlanta:

\begin{lstlisting}[style=sh]
[wx] resolving api.open-meteo.com ...
[wx] api.open-meteo.com = 188.40.99.226
[wx] forecast for 33.749, -84.388  (2026-06-25T20:45 UTC)
[wx]   temperature : 28.5 C
[wx]   wind        : 5.4 km/h from 184 deg
[wx]   conditions  : mainly clear
\end{lstlisting}

That single run exercises the entire chapter: a UDP exchange (DNS) and a TCP
client conversation (HTTP), each dispatched through \texttt{Net\_Devices.%
Device'Class} into the W5500 backend, down through the socket engine and the VDM
transport to the SPI master --- with a receive timeout guarding the wait, and the
W5500's silicon doing the TCP/IP. The application code, though, is just sockets:
the same program, pointed at a desktop \texttt{GNAT.Sockets}, would fetch the same
forecast.

\section{An FTP client --- and the two bugs only the board found}
\label{ch:net:ftp}\index{FTP}\index{flow control}

FTP (RFC 959) is the most demanding client in this chapter, and the most
instructive. It is the first to hold \emph{two} TCP connections at once --- a
persistent \emph{control} connection for commands and replies, and a fresh
\emph{data} connection for each transfer --- and the first to move bulk data in
both directions. \texttt{FTP\_Client} is written, like \texttt{DNS\_Client},
entirely against \texttt{GNAT.Sockets}, so the same source runs on a desktop.

\textbf{Passive mode, always.} An embedded client should only ever make
\emph{outbound} connections, so every transfer uses \texttt{PASV}: the client asks
the server to listen and connects out to it --- no listening socket, friendly to
NAT and firewalls. Two of the W5500's eight sockets are busy during a transfer.

\textbf{Streaming through callbacks.} A file is never buffered whole: downloads
push their bytes to a sink, uploads pull from a source. By the rule from
Chapter~\ref{ch:hal}, both must be library-level and closure-free
(\texttt{No\_Implicit\_Dynamic\_Code}); per-call state travels in a \texttt{Ctx}:

\begin{lstlisting}
type Data_Sink   is access procedure (Ctx : System.Address; Chunk : Byte_Array);
type Data_Source is access procedure
  (Ctx : System.Address; Buf : out Byte_Array; Last : out Natural);

procedure Retrieve (S : in out Session; Path : String;
                    Sink : Data_Sink; Ctx : System.Address; Result : out Status);
procedure Store    (S : in out Session; Path : String;
                    Source : Data_Source; Ctx : System.Address; Result : out Status);
\end{lstlisting}

Errors are a \texttt{Status} (\texttt{Connect\_Failed}, \texttt{Auth\_Failed},
\texttt{Timed\_Out}, \dots); a failure that could leave the control stream out of
sync tears the session down, so \texttt{Is\_Open} stays honest. Two examples drive
it: \texttt{esp32s3\_ftp\_inet} downloads from the public \texttt{ftp.gnu.org}, and
\texttt{esp32s3\_ftp} does an upload round-trip against a local server.

\textbf{The host says it works.} The client is exercised on the build machine
against two real servers --- \texttt{pyftpdlib} and the live \texttt{ftp.gnu.org}
--- doing login, \texttt{SIZE}, \texttt{RETR}, \texttt{STOR} and listings, all
byte-exact. Green across the board.

\textbf{The board disagrees.} Flashed onto the ESP32-S3, the download path worked
first try --- a real \texttt{RETR} from \texttt{ftp.gnu.org} over the W5500. But
uploading a 1\,MiB file failed: the server received 64\,--\,104\,KB --- it varied
between runs --- and then the transfer died. Two bugs were hiding in the W5500 send
path, and \emph{neither could be reproduced on the host}, for the same underlying
reason: a desktop's kernel socket buffers are megabytes deep and its TCP close is
graceful. The W5500's per-socket buffer is \textbf{2\,KB}, and you drive the close
yourself.

\textbf{Bug one: send flow control.} The engine's \texttt{Send} read the free
space in the 2\,KB transmit buffer and, finding it full, returned
\texttt{No\_Space} --- whereupon the upper layer gave up. On a desktop the buffer
never fills (the kernel swallows a megabyte); on the W5500 it fills within the
first hundred kilobytes, the instant the application outruns the wire. The fix is
what TCP flow control is meant to do --- \emph{wait} for the peer to drain the
buffer, rather than abandon the transfer:

\begin{lstlisting}
--  before:  if Free = 0 then Sent := 0; Result := No_Space; return; end if; -- caller quits
--  after:
loop
   Free := R16_Stable (S, Sn_TX_FSR);     --  free space in the TX buffer
   exit when Free > 0;                     --  room now: write and SEND
   if not Connected (S) or else Clock >= Deadline then
      Sent := 0; Result := Timed_Out; return;
   end if;
   Wait_Event (S);                         --  let the peer drain it (it ACKs)
end loop;
\end{lstlisting}

\textbf{Bug two: the missing FIN.} With data now flowing the full megabyte
arrived --- and the transfer \emph{still} hung at the end. After the last byte the
client closes the data connection to signal end-of-file; the server reads to EOF,
then replies \texttt{226 Transfer complete}. But \texttt{Close} was issuing the
W5500's abrupt \texttt{CLOSE} command, which discards the socket without sending a
\texttt{FIN}. The server never saw end-of-stream, never sent its \texttt{226}, and
the client timed out waiting. On a desktop the OS sends the \texttt{FIN} for you;
here you must ask for it --- an active close (\texttt{DISCON}, which sends
\texttt{FIN} and waits for the handshake) before the socket is freed.

With both fixes the upload completes, and the example reads the file back and
checks it byte-for-byte:

\begin{lstlisting}[style=sh]
[ftp] STOR /from_board.bin (1048576 bytes): OK
[ftp] read-back 1048576 bytes: round-trip VERIFIED
\end{lstlisting}

\textbf{The lesson.} Both bugs lived in code that had passed every host test and a
careful reading --- including an analysis that ``proved,'' from the source, that
the transfer would sustain. Neither was a logic error visible on the page: each was
a property of \emph{real hardware under real load} --- a small buffer that actually
fills, a close you actually have to drive. The portable layer let the protocol be
written and tested in minutes on a laptop; the silicon is what proved it. On bare
metal, the board is the test.

\section{An FTP server, and an accessibility check that earned its keep}
\label{ch:net:ftpd}\index{FTP!server}\index{accessibility check}

The other half of FTP is the server, and it joins the two halves of this book: it
serves the \emph{ext4-on-flash} filesystem of Chapter~\ref{ch:storage} over the
network --- and, through a small mount table, any other ext4 volume alongside it in
one tree.

\subsection{One namespace for several volumes: the VFS}
\label{ch:net:vfs}\index{VFS}\index{mount point}

A board may have more than one place to keep files --- the on-board SPI flash and
an SD card, say. Both already reduce to the same abstraction: each is a
\texttt{Block\_Dev.Device} (the flash through the wear-levelling FTL, the card
through \texttt{Block\_Dev.SDMMC\_Source}), and each mounts as an
\texttt{ESP32S3.Ext4.FS.Mount}. What is left is to \emph{name} them and route a path
to the right one. That is all \texttt{ESP32S3.Ext4.VFS} does --- a tiny mount table:

\begin{lstlisting}
ESP32S3.Ext4.VFS.Add ("flash", Flash_M'Access);   -- appears at /flash
ESP32S3.Ext4.VFS.Add ("sd",    SD_M'Access);       -- appears at /sd   (optional)
\end{lstlisting}

Each \texttt{Add} registers a mounted filesystem under a top-level name. A path is
then resolved by its first component: \texttt{Resolve ("/flash/docs/a.txt")} returns
the flash mount and the remainder \texttt{"/docs/a.txt"}; the bare root \texttt{"/"}
is \emph{virtual} --- it belongs to no device and simply lists the registered names.
This is the Unix mount-point model in miniature, and the storage analogue of the
chip-neutral NIC registry of \S\ref{ch:net:netdevice}: a neutral device interface
(\texttt{Block\_Dev.Device}) plus a small table mapping a name to a backend. With
one device registered the tree is just \texttt{/flash}; adding the SD card later is
the single \texttt{Add} line above --- nothing in the server changes.

\texttt{FTP\_Server.Run} therefore carries no filesystem of its own; it serves
whatever has been registered:

\begin{lstlisting}
ESP32S3.Ext4.VFS.Add ("flash", Flash_FS.M'Access);
FTP_Server.Run (Server_Name => "Ada ESP32-S3");   -- the 220 greeting
\end{lstlisting}

Note what is \emph{absent}: no address. Passive FTP must tell the client which IP
to open the data connection to, and rather than hand-feed it the DHCP lease, the
server reads its own address from each accepted control connection
(\texttt{Get\_Socket\_Name} returns the interface's configured IP) and advertises
exactly that in the \texttt{227} reply. So it always names the address the client
actually reached the board on, even as DHCP reassigns it. \texttt{Run} does take an
optional \texttt{Local\_IP} to override this --- needed only behind NAT, where the
public address differs from the board's own --- but ordinary use leaves it out.

It runs the accept loop --- one client at a time, a control connection plus one
passive data connection per transfer. The server side of the socket API is
\texttt{Bind\_Socket} / \texttt{Listen\_Socket} / \texttt{Accept\_Socket} (the echo
example, \S\ref{ch:net:examples}). Each command first resolves its path through the
VFS, then maps onto an ext4 operation on the chosen mount --- \texttt{LIST} onto
\texttt{Iterate} (plus \texttt{Stat} for the \texttt{ls -l} columns), \texttt{RETR}
onto the offset \texttt{Read\_File}, \texttt{STOR} onto \texttt{Create\_File} +
\texttt{Append} + \texttt{Commit}, \texttt{DELE} / \texttt{MKD} / \texttt{RMD} onto
\texttt{Unlink} / \texttt{Mkdir} / \texttt{Rmdir} --- while \texttt{LIST} of the
virtual root just enumerates the mount names. The filesystem's
\texttt{Ada.IO\_Exceptions} are caught and turned into FTP \texttt{550} replies, so
a bad request never takes the server down; \texttt{Read\_Only} refuses the mutating
commands. A desktop client --- \texttt{ftp}, FileZilla, a browser, Python
\texttt{ftplib} --- then sees \texttt{/flash} (and \texttt{/sd}, when present) as
directories at the root, and can browse, download, upload and manage the board's
storage beneath them.

\textbf{The board faulted on the first run --- and that was the language working.}
Flashed and booted, the server printed its greeting and then died:

\begin{lstlisting}[style=sh]
[ftpd] serving on 192.168.1.229:21  (anonymous, read-write)
Unhandled Ada Exception: PROGRAM_ERROR
Message: ftp_server.adb:539 accessibility check failed
\end{lstlisting}

Two things are worth drawing out. First, the \emph{diagnosis}: from the host the
board still answered \texttt{ping}, yet port 21 refused connections. That
combination is a tell --- the W5500 answers ARP and ICMP \emph{in silicon} once its
address is configured, so a board that pings but refuses every port has a
\emph{halted CPU}, not a missing listener. The chip was up; the program was not.

Second, the fault itself. Line 539 was the server stashing the mount for its
handlers to reach:

\begin{lstlisting}
Mnt := FS;   -- Mnt is library-level; FS points at the caller's Mount
\end{lstlisting}

The example had declared the \texttt{Mount} as a local variable of \texttt{Main}.
Storing an access to it in a library-level variable --- one that outlives
\texttt{Main}'s frame --- is exactly the use-after-scope that Ada's
\emph{accessibility rules} forbid, and the one case the compiler cannot prove
statically it checks at run time. The check fired the instant the assignment ran,
turning a latent dangling pointer into a located exception with a line number. The
fix is to give the filesystem objects the lifetime the server assumes --- declare
them at library level:

\begin{lstlisting}
package Flash_FS is                       -- library level: lives for the program
   Raw : aliased ESP32S3.Block_Dev.W25Q_Source.Source;
   Vol : aliased ESP32S3.Block_Dev.WL.Volume;
   Dev : ESP32S3.Block_Dev.Device;
   M   : aliased ESP32S3.Ext4.FS.Mount;
end Flash_FS;
\end{lstlisting}

The \texttt{Mnt := FS} store shown above belonged to the original single-mount
server; the move to the VFS of \S\ref{ch:net:vfs} did not retire the rule, only
relocated it. \texttt{VFS.Add} keeps each mount's access in its table, so a
registered filesystem must outlive the server just the same --- which is exactly
why the example passes \texttt{Flash\_FS.M'Access}, a library-level object, and not
a \texttt{Mount} local to \texttt{Main}. The accessibility check guards the new
shape as faithfully as it guarded the old.

In C the same mistake is a silent use-after-scope that corrupts something later, far
from the cause. Here it was a clean exception at the assignment that made it. The
two FTP halves bracket the lesson nicely: the client's bugs were properties of the
hardware that no amount of reading the source could reveal; the server's was a
property of the \emph{program} that the language caught before it could do harm.

\noindent\textbf{Example:} \texttt{examples/esp32s3\_ftp\_server} --- formats a
fresh ext4 on the W25Q flash, registers it at \texttt{/flash}, then serves it;
HW-verified against \texttt{ftplib} (\texttt{NLST "/"} returns \texttt{flash};
\texttt{LIST}, \texttt{RETR}/\texttt{SIZE}, \texttt{STOR} with a byte-exact
read-back, \texttt{MKD}, \texttt{DELE} all under \texttt{/flash}). An SD card would
join the tree at \texttt{/sd} with one more \texttt{VFS.Add} call and no change to
the server.

\section{Multiple interfaces: routing, metrics, and pinning}
\label{ch:net:multinic}\index{routing}\index{interface!multiple}\index{failover}

A board may carry more than one way onto a network --- two W5500s, or a wired NIC
plus a cellular modem. The registry of \S\ref{ch:net:netdevice} already allows it:
\texttt{Add\_Interface} registers several \texttt{Net\_Devices.Device}s, and a
\texttt{Socket\_Type} records which one it is on. What that section left open ---
its own comment said routing was \emph{``not wired yet''} --- is the policy: which
interface a given connection should use. This section is that policy.

It rests on one new fact each interface reports, \texttt{Is\_Up}: is it usable right
now --- physically up, with an address? (For the W5500, PHY link up and a non-zero
IP.) Routing consults it so traffic only leaves a live interface; pinning uses it to
fail closed.

\subsection{The routing table}
\index{Net\_Routes}\index{metric}

\texttt{Net\_Routes} is a small table --- the network analogue of the VFS mount
table of \S\ref{ch:net:vfs}, a neutral interface plus a table that maps a key (here
a destination address) to a backend (here an interface):

\begin{lstlisting}
Net_Routes.Set_Default (Eth,  Metric => 10);    -- 0.0.0.0/0 out the wired NIC
Net_Routes.Set_Default (Cell, Metric => 100);   -- ... or cellular, less preferred
Net_Routes.Add_Route ((10,0,0,0), (255,0,0,0), Eth, Metric => 10);  -- 10/8 is wired
\end{lstlisting}

A destination is resolved in a fixed order: among routes that \emph{match} and whose
interface is \emph{up}, take the \textbf{longest prefix}, then the \textbf{lowest
metric}. The two keys do different jobs. Prefix length is specificity --- a
\texttt{10.0.0.0/8} route beats a \texttt{0.0.0.0/0} default regardless of metric, so
``LAN traffic goes out the wired port'' is expressed by a more specific route, not a
number. Metric breaks ties \emph{among equally-specific routes} --- two defaults, one
per interface, are decided by metric, so the wired default (10) wins over the
cellular default (100). ``Cellular is billed --- always last'' is a metric.

The \texttt{Is\_Up} filter is what turns preference into \textbf{failover}. Give both
interfaces a default route, wired cheaper; normally the wired default wins. Unplug the
wired link and its routes drop out of consideration, so the cellular default takes
over --- automatically, no application code. Plug it back in and it preempts again
(lower metric, and now up). One subtlety worth stating: an existing connection stays
on the interface it was opened on (the socket remembers it); failover re-routes only
\emph{new} connections. You cannot migrate a live TCP socket between interfaces, and
the design does not pretend to.

The table is pure logic --- liveness is \emph{injected} (the facade hands it a query
backed by the registry), so it is unit-tested natively with a mock up-state, no
hardware: longest-prefix, metric tie-break, failover, preempt-back, and the no-route
cases all checked on the host before any of it touches silicon.

\subsection{Pinning: this link and no other}
\index{pinning}\index{interface!pinning}

Routing is for ``reach this sensibly, survive a link loss.'' The opposite need ---
``this traffic is allowed on exactly one link and nowhere else'' --- is
\textbf{pinning}:

\begin{lstlisting}
GNAT.Sockets.Set_Interface (S, Eth);   -- S uses Eth only
\end{lstlisting}

A pinned socket bypasses the table entirely: \texttt{Connect} uses that interface and
\textbf{fails closed} --- \texttt{Socket\_Error} if it is down --- rather than falling
back. That difference is the whole point. ``Prefer wired'' (a metric) silently fails
over; ``only wired'' (a pin) must not, or it is a security or billing bug. Binding a
socket to a specific interface's own address pins it the same way, so a server can be
made to listen on just one interface. Inbound is otherwise pinned for free --- a
listener lives on the chip it was created on.

\subsection{How a connect chooses}

For an \emph{unpinned} socket, \texttt{Connect\_Socket} resolves the destination
through the table and then re-homes the socket onto the chosen interface before it
opens the chip socket. The design stays invisible to single-interface boards by
construction: with no routes configured the resolver returns the default interface
and the re-home is a no-op, so every existing client --- the weather fetch, the FTP
client, DNS --- behaves exactly as before. Only once routes exist does a destination
with no live route raise \texttt{Socket\_Error} instead of leaking out some default.

\noindent\textbf{Example:} \texttt{examples/esp32s3\_multinic} brings up two W5500s
--- the primary by DHCP, a second on \texttt{CS=IO40} whose reset is shared on
\texttt{IO11} (so it uses a software reset, \texttt{MR.RST}, leaving the primary
alone) and which is left uncabled on purpose. The board reports \texttt{eth0} up and
\texttt{eth1} down; routes \texttt{8.8.8.8} out \texttt{eth0}; falls a \texttt{10/8}
destination back to the default when its specific route (\texttt{eth1}) is down;
\emph{refuses} a connect pinned to the down \texttt{eth1}; and connects for real out
the live \texttt{eth0} --- the routing table, liveness, and fail-closed pinning, all
on hardware.

\subsection{Beyond W5500: cellular and friends}

The reason a cellular modem slots in here is that \texttt{Net\_Devices.Device} is a
\emph{socket-level} contract (\texttt{Connect}/\texttt{Listen}/\texttt{Send} by
address and port), not a packet-level one. A hardware-offload NIC like the W5500 and a
firmware-stack modem driven by AT commands both implement it directly, with no on-chip
TCP/IP stack --- the modem's \texttt{Connect} issues an AT command, the W5500's writes
a socket register, and the facade above neither knows nor cares. A new interface ---
another W5500, or \texttt{Cell\_Modem.Net\_Device} --- is one more
\texttt{Add\_Interface} call and a route or two; nothing above the facade changes. (A
raw Ethernet MAC or Wi-Fi part would \emph{not} fit so cleanly: those are packet-level
and would need a software IP stack beneath this interface.)
